\section{Discussion, prior art, and open directions}
\label{sec:discussion}

\subsection{The rational--integral contrast}

The main result has a deliberately asymmetric scope.  The minimum regulus
trades generate every rational relation in all three configurations.  At the
integral level we prove saturation at
\(2,3,17,137,219\,097,3\,288\,036\,131\).  Globally, the exact
two-minor certificate closes one of the six large maps, namely the complete
\(\rh/\cB_1\) component, away from \(2,3,17\).  Equality of the full
trade and relation lattices remains open because the other five large
component gates have not been computed.  Thus
\[
\text{minimum local geometry}
\quad\Longrightarrow\quad
\text{full rational generation}
\ \centernot\Longrightarrow\
\text{global integral saturation}.
\]
The character table also shows that common incidence parameters and a
common group do not make the relation modules isomorphic.

The general theories used here have distinct ownership.  Cho's
characteristic-zero null-design decomposition and minimum subspace
null-design work provide the complete-Grassmann background
\cite{Cho1997,Cho1998}.  Krotov supplies the minimum-support theorem and
the \(t=1\) classification \cite{Krotov2017}.  The ordinary-subset
orbit-spanning result of Ghorbani--Kamali--Khosrovshahi is the closest
structural analogue of \cref{thm:full-kernel}
\cite{GhorbaniKamaliKhosrovshahi}.  It uses the full symmetric action on
subsets, whereas the present theorem concerns selected binary subspace
constituents and the order-\(204\) group.

For toric ideals, Aoki--Takemura--Yoshida own the indispensable-fiber
criterion \cite{AokiTakemuraYoshida}.  Betti--Grate--Holleben--Salizzoni
develop the current incidence-toric/null-design language for ordinary
subset incidence matrices \cite{BettiEtAl2026}.  The present finite
calculation is the proof that the cubic fibers of these selected
Grassmann configurations have exactly the two regulus monomials.
Classical integral-design modules and circuit lattices are represented by
\cite{GraverJurkat,GrahamLiLi}; the component computation in
\cref{sec:137} is specific to the locked triple.  These comparisons are
scope boundaries, not priority claims.

Fitting ideals, cyclotomic arithmetic, CRT reconstruction, and Smith normal
form are standard tools \cite{StacksProjectFitting,Cohen1993,Washington1997,Newman1972}.
The contribution here is their exact application to the frozen complete
component map and the independently replayed two-minor certificate.  No
method-firstness or priority claim is made.

\subsection{Tomography and the P42 interface}

There is a useful but logically separate sensing result.  The
\(97\,155\) three-dimensional subspaces of \(\FF_2^8\) split into
\(515\) \(\GammaG\)-orbits, with orbit-size profile
\[
 51^5,\qquad102^{70},\qquad204^{440}.
\]
The selected three-space versus four-space incidence map
\[
 W_{3,4}:\QQ[\cB_i]\longrightarrow\QQ[\operatorname{Gr}_2(8,3)]
\]
has column rank \(66\,929\), so it is injective on the entire block
module and, in particular, on every \(K_i\).  It gives a lossless rational
sensor system for the relations.

This statement is not used in
\cref{thm:full-kernel,thm:rho137,thm:theta137,thm:theta219097,thm:theta3288036131,thm:global-rho}.
The temporary selected-row rank losses at \(137\), \(219\,097\), and
\(3\,288\,036\,131\) were
presentation phenomena, not tomography failures and not integral quotient
classes.  Conversely, no tomography statement is used in the exact
cyclotomic or Fitting argument of \cref{thm:global-rho}.  The broader
comparison of sensor families,
minimal sensor sets, and reconstruction belongs to the planned P42
tomography paper.  The synchronized interface is:
\[
\begin{array}{c|c}
\text{P39}&\text{P42}\\
\midrule
\text{relation modules and circuit generators}
&\text{measurement maps and reconstruction}\\
\text{one lossless \(W_{3,4}\) check}
&\text{systematic sensor comparison}\\
\text{integral component audit}
&\text{distinguish sensing loss from lifting loss}.
\end{array}
\]

\subsection{Remaining global integral gates}

At \(137\), \(219\,097\), and \(3\,288\,036\,131\), all six large maps
are classified and
\cref{thm:theta137,thm:theta219097,thm:theta3288036131} gives complete
three-design vanishing.  These were exactly the good-prime signals found
by the frozen selected-row search.  Their disappearance under complete-map
audits shows that the candidate list was not a support theorem.

\Cref{thm:global-rho} goes beyond that finite list for one map: every
prime ideal outside \(2,3,17\) is excluded from the complete
\(\rh/\cB_1\) cokernel.  The remaining integral tasks are therefore:
\begin{itemize}
  \item compute an exact localized maximal-minor ideal certificate for each
  of the other five large component maps;
  \item combine those results with \cref{prop:small-components} to prove or
  refute \(\operatorname{Supp}(\Theta_i)\subseteq\{2,3,17\}\);
  \item compute the exact global Smith invariant factors, including the
  remaining \(2\)-, \(3\)-, and \(17\)-primary structure.
\end{itemize}
A further finite-prime search cannot close these questions.  The successful
\(\rh/\cB_1\) route shows the required form of a certificate: one common
integral cyclotomic basis, exact characteristic-zero determinants from
genuinely different row bases, and their ideal sum rather than a rational
norm gcd.  The failure of closely related same-row minors in
\cref{rem:minor-artifact} is the reason row-basis diversity matters.
The five repetitions are deliberately deferred; none is needed for the
rational, cubic, selected-prime, or one-component global theorems proved
here.
\subsection{Geometric and family questions}

Krotov's bound gives support \(30\) and leg volume \(15\) for
\(T_2(2,4,8)\), but the \(t=1\) uniqueness theorem used in this paper does
not classify that case.  An exact next question is a two-colour census:
one leg in \(\cB_i\), the other in \(\cB_j\).  The canonical construction
can be searched orbitally, but exhaustiveness requires a separate
classification theorem.

A second direction is to replace the locked finite triple by a general
selected \(q\)-subspace configuration.  The missing theorem is not the
minimum-trade bound; it is nonzero projection of the selected minimum
trades to every null-design constituent.  The component-rank criterion in
\cref{lem:component-rank} isolates exactly what such a theorem would have
to prove.

Finally, the higher-degree parts of \(I_i\) remain open.  Cubic
indispensability neither implies cubic generation nor determines the
complete Markov or Graver basis.  The P39 and P42 programs meet naturally
at the sequence
\[
 \QQ\text{-generation}\longrightarrow
 \ZZ\text{-lifting}\longrightarrow
 \text{Markov connectivity},
\]
but the three arrows should remain separate until their respective
integral and higher-degree obstructions are classified.
